\chapter{โครงงานย่อย การสังเคราะห์และสาธิตหลักการฟิสิกส์ในเครื่องมือห้องปฏิบัติการจุลชีววิทยา}

\begin{planbox}[สัปดาห์ที่ 14 — เวลาปฏิบัติการ 3 ชั่วโมง (สอดคล้องกับ CLO2, CLO3 และ CLO4 — คะแนนเต็ม 15 คะแนน)]
\textbf{1. วัตถุประสงค์ของโครงงานย่อย (Mini-Project)}
\begin{itemize}
    \item สังเคราะห์องค์ความรู้ทางฟิสิกส์ตลอดทั้งภาคการศึกษา (กลศาสตร์ ความร้อน อุณหพลศาสตร์ ไฟฟ้า และทัศนศาสตร์) เพื่ออธิบายกลไกการทำงานของเครื่องมือวิทยาศาสตร์ในห้องปฏิบัติการจุลชีววิทยาอย่างน้อย 1 เครื่องอย่างลึกซึ้ง
    \item สาธิตขั้นตอนการใช้งานเครื่องมือ การคำนวณพารามิเตอร์ตั้งค่าที่ถูกต้อง และการป้องกันอุบัติเหตุตามหลักความปลอดภัยห้องปฏิบัติการ
    \item นำเสนอผลงานในรูปแบบโปสเตอร์อินโฟกราฟิกหรือการสาธิตหน้าชั้นเรียน (10 นาทีต่อกลุ่ม) และตอบข้อซักถามเชิงแก้ปัญหา (Troubleshooting) ได้อย่างถูกต้อง
\end{itemize}
\textbf{2. ความเชื่อมโยงกับผลลัพธ์การเรียนรู้ของรายวิชา (CLOs) และหลักสูตร (PLOs)}
\begin{itemize}
    \item \textbf{CLO2 (PLO2, PLO3)}\quad ประยุกต์ใช้หลักการทางฟิสิกส์อธิบายกลไกและแก้ไขปัญหาการใช้เครื่องมือวิทยาศาสตร์ในงานจุลชีววิทยาได้ (5 คะแนน)
    \item \textbf{CLO3 (PLO1)}\quad สาธิตมาตรการความปลอดภัยและจรรยาบรรณในการดูแลรักษาเครื่องมือได้อย่างครบถ้วน (5 คะแนน)
    \item \textbf{CLO4 (PLO5, PLO6)}\quad ทำงานร่วมกันเป็นทีมและสื่อสารข้อมูลวิทยาศาสตร์ผ่านโปสเตอร์และการนำเสนอได้อย่างชัดเจน (5 คะแนน)
\end{itemize}
\end{planbox}

\section{แนวทางการเลือกหัวข้อโครงงานย่อยและกรอบการวิเคราะห์เชิงฟิสิกส์}

นักศึกษาแต่ละกลุ่ม (กลุ่มละ 3--4 คน) เลือกเครื่องมือหรือเทคนิคปฏิบัติการทางจุลชีววิทยา 1 หัวข้อจากรายการที่กำหนดให้ หรือเสนอหัวข้อเพิ่มเติมต่อผู้สอน โดยต้องวิเคราะห์เชื่อมโยงให้ครบ 4 มิติ ได้แก่ \textbf{(1) หลักการฟิสิกส์พื้นฐานและสมการหลัก} \textbf{(2) ตัวอย่างการคำนวณตั้งค่าใช้งานจริง} \textbf{(3) สาเหตุความผิดพลาดที่พบบ่อยและวิธีแก้ไข} และ \textbf{(4) กฎความปลอดภัยและจรรยาบรรณการบำรุงรักษา}

\begin{conceptbox}[หัวข้อตัวอย่างสำหรับการเลือกทำโครงงานย่อย (เลือกกลุ่มละ 1 หัวข้อ)]
\begin{enumerate}
    \item \textbf{หัวข้อที่ 1 ฟิสิกส์ของหม้อนึ่งความดันไอ (Autoclave)}\quad ความดันไออิ่มตัว กฎของแก๊ส ความร้อนแฝงของการควบแน่นไอน้ำ และการตรวจสอบความสำเร็จของการฆ่าเชื้อสปอร์แบคทีเรียที่ $121\,^\circ\text{C}$ ($15\text{ psi}$)
    \item \textbf{หัวข้อที่ 2 ฟิสิกส์ของเครื่องปั่นเหวี่ยง (Centrifuge)}\quad แรงสู่ศูนย์กลาง กฎของสโตกส์ การแปลงค่าระหว่าง $\text{RPM}$ กับ $\text{RCF}$ ($\times g$) และฟิสิกส์ของการถ่วงสมดุลโรเตอร์
    \item \textbf{หัวข้อที่ 3 ฟิสิกส์ของกล้องจุลทรรศน์และน้ำมันช่วยส่อง (Microscopy \& Immersion Oil)}\quad ทางเดินแสงในเลนส์ประกอบ สมการการแยกชัดของเอบเบอ ($d_{\min}$) และการจับคู่ดัชนีหักเหที่หัวเลนส์ $100\times$
    \item \textbf{หัวข้อที่ 4 ฟิสิกส์ของเครื่องสเปกโตรโฟโตมิเตอร์และการวัด $\text{OD}_{600}$}\quad กฎของเบียร์--แลมเบิร์ต การกระเจิงแสงของเซลล์ และการแก้ไขข้อจำกัดเมื่อความขุ่นเกินช่วงเชิงเส้น
    \item \textbf{หัวข้อที่ 5 ฟิสิกส์ของตู้ปลอดเชื้อ (Biosafety Cabinet)}\quad พลศาสตร์การไหลแบบราบเรียบ (Laminar Airflow) แผ่นกรอง HEPA และการลดลงของความเข้มรังสี UV-C ตามกฎกำลังสองผกผัน
    \item \textbf{หัวข้อที่ 6 ฟิสิกส์ของชุดแยกสารพันธุกรรมด้วยไฟฟ้า (Gel Electrophoresis)}\quad กฎของโอห์ม สนามไฟฟ้า กำลังความร้อนจูลในบัฟเฟอร์ และแรงหนืดต้านการเคลื่อนที่ของดีเอ็นเอในเจล
\end{enumerate}
\end{conceptbox}

\begin{figure}[H]
\centering
\resizebox{0.86\textwidth}{!}{%
\begin{tikzpicture}[>=Stealth, font=\small]
    \node[draw=physRoyalBlue, fill=boxBlueBg, line width=1.2pt, rounded corners=6pt, minimum width=3.2cm, minimum height=1.8cm, align=center, font=\bfseries] (s1) at (0,0) {ขั้นที่ 1\\เลือกเครื่องมือจุลชีววิทยา\\และระบุหลักฟิสิกส์};
    \node[draw=bioEmerald, fill=boxGreenBg, line width=1.2pt, rounded corners=6pt, minimum width=3.2cm, minimum height=1.8cm, align=center, font=\bfseries] (s2) at (4.2,0) {ขั้นที่ 2\\แสดงตัวอย่างคำนวณจริง\\(แบบ Low-Math)};
    \node[draw=crimsonAlert, fill=boxRedBg, line width=1.2pt, rounded corners=6pt, minimum width=3.2cm, minimum height=1.8cm, align=center, font=\bfseries] (s3) at (8.4,0) {ขั้นที่ 3\\วิเคราะห์ข้อผิดพลาด\\และกฎความปลอดภัย};
    \node[draw=physAmber!90!black, fill=boxAmberBg, line width=1.2pt, rounded corners=6pt, minimum width=3.2cm, minimum height=1.8cm, align=center, font=\bfseries] (s4) at (12.6,0) {ขั้นที่ 4\\จัดทำโปสเตอร์และสาธิต\\หน้าชั้นเรียน (15 คะแนน)};
    \draw[->, line width=1.5pt, physRoyalBlue] (s1) -- (s2);
    \draw[->, line width=1.5pt, bioEmerald] (s2) -- (s3);
    \draw[->, line width=1.5pt, crimsonAlert] (s3) -- (s4);
\end{tikzpicture}%
}
\caption{ขั้นตอนการดำเนินโครงงานย่อย (Mini-Project) ฟิสิกส์ในเครื่องมือห้องปฏิบัติการจุลชีววิทยาในสัปดาห์ที่ 14}
\label{fig:lab13_workflow}
\end{figure}

\section{แบบฟอร์มสรุปโครงงานย่อยและใบประเมินคะแนน (15 คะแนน)}

\begin{worksheetbox}[แบบฟอร์มสรุปโครงงานย่อยและบันทึกคะแนนการนำเสนอ (สัปดาห์ที่ 14)]
\textbf{กลุ่มปฏิบัติการที่} \dotfill \textbf{วันที่นำเสนอ} \dotfill\\[4pt]
\textbf{ชื่อหัวข้อโครงงานย่อยและเครื่องมือที่ศึกษา} \dotfill\\[4pt]
\textbf{รายชื่อสมาชิกและหน้าที่รับผิดชอบ}\\[2pt]
1. \dotfill หน้าที่ \dotfill\\[2pt]
2. \dotfill หน้าที่ \dotfill\\[2pt]
3. \dotfill หน้าที่ \dotfill\\[2pt]
4. \dotfill หน้าที่ \dotfill

\vspace{6pt}
\textbf{1. สรุปหลักการทางฟิสิกส์และสมการหลักที่ควบคุมการทำงานของเครื่องมือนี้ (CLO2)}\\[4pt]
\dotfill\\[6pt]
\dotfill\\[6pt]
\dotfill

\vspace{6pt}
\textbf{2. ตัวอย่างการคำนวณพารามิเตอร์ในการใช้งานจริง 1 ตัวอย่างพร้อมหน่วยวัด SI}\\[4pt]
\dotfill\\[6pt]
\dotfill\\[6pt]
\dotfill

\vspace{6pt}
\textbf{3. แนวทางการแก้ปัญหาเมื่อเกิดข้อผิดพลาด (Troubleshooting) และกฎความปลอดภัย (CLO3)}\\[4pt]
\dotfill\\[6pt]
\dotfill

\vspace{6pt}
\textbf{ตารางที่ 13.1\quad ตารางบันทึกคะแนนประเมินโครงงานย่อยสำหรับอาจารย์ผู้สอน (เต็ม 15 คะแนน)}
\begin{center}
\footnotesize\setlength{\tabcolsep}{3.5pt}
\begin{tabularx}{\linewidth}{>{\raggedright\arraybackslash}X c c c}
\toprule
\textbf{รายการประเมินตามผลลัพธ์การเรียนรู้ (OBE Rubrics)} & \textbf{CLOs ที่วัด} & \textbf{คะแนนเต็ม} & \textbf{คะแนนที่ได้} \\
\midrule
1. ความถูกต้องของการอธิบายกลไกทางฟิสิกส์ การคำนวณ และการแก้ปัญหาหน้างาน & CLO2 & $5\text{ คะแนน}$ & \dotfill \\[6pt]
2. การสาธิตขั้นตอนความปลอดภัยในการใช้เครื่องมือและจรรยาบรรณห้องปฏิบัติการ & CLO3 & $5\text{ คะแนน}$ & \dotfill \\[6pt]
3. การทำงานร่วมกันของสมาชิกในทีม ความชัดเจนของโปสเตอร์ และการสื่อสารตอบคำถาม & CLO4 & $5\text{ คะแนน}$ & \dotfill \\
\midrule
\textbf{รวมคะแนนโครงงานย่อยประจำกลุ่ม (คิดเป็น 15\% ของคะแนนรวมรายวิชา)} & \textbf{CLO2--4} & \textbf{15 คะแนน} & \dotfill \\
\bottomrule
\end{tabularx}
\end{center}

\vspace{4pt}
\textbf{ข้อเสนอแนะจากอาจารย์ผู้สอน} \dotfill \textbf{ลงชื่อผู้ประเมิน} \dotfill
\end{worksheetbox}
